\documentclass[11pt,a4paper]{article}

% ── Packages ────────────────────────────────────────────────────────────────
\usepackage[utf8]{inputenc}
\usepackage[T1]{fontenc}
\usepackage{lmodern}
\usepackage[margin=2.5cm]{geometry}
\usepackage{amsmath,amssymb}
\usepackage{graphicx}
\usepackage{booktabs}
\usepackage{hyperref}
\usepackage{natbib}
\usepackage{xcolor}
\usepackage{microtype}
\usepackage{siunitx}
\usepackage{caption}
\usepackage{subcaption}
\usepackage{multirow}
\usepackage{lineno}

% ── Metadata ────────────────────────────────────────────────────────────────
\hypersetup{
  pdftitle={Climate-Driven Energy Demand Shifts in the Kisalfold Region},
  pdfauthor={Author Name},
  colorlinks=true,
  linkcolor=blue!60!black,
  citecolor=blue!60!black,
  urlcolor=blue!60!black,
}

% ── Custom commands ──────────────────────────────────────────────────────────
\newcommand{\HDD}{\mathrm{HDD}}
\newcommand{\CDD}{\mathrm{CDD}}
\newcommand{\degC}{$^\circ$C}
\newcommand{\todo}[1]{\textcolor{red}{\textbf{[TODO: #1]}}}

% ── Line numbers (remove before submission) ──────────────────────────────────
\linenumbers

% ============================================================================
\begin{document}

\title{Climate-Driven Energy Demand Shifts in the Kisalf\"{o}ld Region:\\
Historical Correlations, Future Projections,\\
and Infrastructure Loss Implications}

\author{
  Author Name\thanks{Corresponding author. \texttt{your@email.com}}\\
  \small P\'{a}pa, Kisalf\"{o}ld, Hungary
}

\date{\today}

\maketitle

% ── Abstract ────────────────────────────────────────────────────────────────
\begin{abstract}
Rising temperatures driven by anthropogenic climate change are expected to 
substantially alter patterns of energy consumption across Central Europe. 
This study investigates the historical relationship between surface temperature 
and electricity demand in the Kisalf\"{o}ld region of Hungary --- a lowland 
agricultural area characterised by high summer temperature variability --- 
and projects how this relationship evolves under CMIP6 warming scenarios 
of \SI{1.5}{\celsius}, \SI{2}{\celsius}, and \SI{3}{\celsius} above 
pre-industrial levels by 2050. Using ERA5 reanalysis data spanning 1970--2024, 
local twice-daily atmospheric soundings, and Hungarian national grid load 
records from MAVIR and ENTSO-E, we train an XGBoost demand model achieving 
\todo{RMSE} on held-out test years. Projections suggest 
\todo{X--Y}\% increases in annual cooling degree days by mid-century under 
a \SI{2}{\celsius} scenario. Critically, we estimate that 
\todo{Z}\% of the associated increase in generated electricity will be 
lost to grid transmission inefficiencies exacerbated by elevated ambient 
temperatures, representing a compounded infrastructure challenge. 
An interactive Progressive Web Application accompanying this paper makes 
all projections publicly accessible.
\end{abstract}

\textbf{Keywords:} climate change, energy demand, Kisalf\"{o}ld, Heating 
Degree Days, Cooling Degree Days, XGBoost, CMIP6, grid transmission losses, 
Hungary

\hrule
\vspace{1em}

% ── 1. Introduction ──────────────────────────────────────────────────────────
\section{Introduction}
\label{sec:introduction}

The relationship between ambient temperature and electricity demand is 
well-established in the literature \citep{\todo{ref}}. As global mean 
temperatures rise, two opposing forces reshape regional demand profiles: 
declining heating requirements in winter months and sharply increasing 
cooling requirements in summer, driven by growing air conditioning 
adoption \citep{\todo{ref}}. The net effect is non-linear, 
geographically heterogeneous, and sensitive to the socioeconomic context 
of the affected region.

The Kisalf\"{o}ld (\textit{Little Hungarian Plain}) is a lowland region 
straddling northwestern Hungary and eastern Austria, characterised by 
continental climate dynamics, intensive agriculture, and increasing 
summer heat stress. \todo{Add 1--2 sentences on why this region is 
particularly interesting/vulnerable.}

This study makes three contributions:
\begin{enumerate}
  \item A historical correlation analysis of temperature and grid load 
    for Hungary using ERA5 and ENTSO-E data (1970--2024), cross-validated 
    against locally collected atmospheric soundings from P\'{a}pa (2025--present).
  \item Demand projections under three CMIP6 warming scenarios to 2050, 
    using a gradient-boosted regression model with SHAP-based interpretability.
  \item A novel quantification of the compounded infrastructure loss effect --- 
    how heat-induced increases in transmission losses amplify the gap between 
    generated and consumed electricity.
\end{enumerate}

% ── 2. Data ──────────────────────────────────────────────────────────────────
\section{Data}
\label{sec:data}

\subsection{ERA5 Reanalysis}
\label{sec:data_era5}

Surface temperature (\SI{2}{\metre} air temperature), solar radiation 
(surface downwelling shortwave), and wind components were obtained from 
the ERA5 reanalysis \citep{hersbach2020era5} via the Copernicus Climate 
Data Store at \SI{0.25}{\degree} horizontal resolution. Data were spatially 
averaged over the Kisalf\"{o}ld bounding box 
($\SI{47.2}{\degree}$--$\SI{47.8}{\degree}$N, 
$\SI{17.0}{\degree}$--$\SI{18.0}{\degree}$E) 
for the period 1970--2024.

\subsection{Local Atmospheric Soundings}
\label{sec:data_soundings}

\todo{Describe your sounding collection: location, frequency, instrument, 
period. Mention validation against OMSZ Gy\H{o}r station.}

\subsection{Hungarian Grid Load Data}
\label{sec:data_load}

Hourly national electricity load data for Hungary were obtained from 
the ENTSO-E Transparency Platform \citep{entsoe} 
(bidding zone \texttt{10YHU-MAVIR----U}) for 2015--2024, 
and from MAVIR annual reports for earlier years. \todo{Describe any 
preprocessing, missing data handling.}

\subsection{CMIP6 Climate Projections}
\label{sec:data_cmip6}

Future temperature trajectories were derived from the CMIP6 multi-model 
ensemble \citep{eyring2016cmip6} under Shared Socioeconomic Pathways 
SSP1-2.6 (\SI{1.5}{\celsius}), SSP2-4.5 (\SI{2}{\celsius}), and 
SSP5-8.5 (\SI{3}{\celsius}). \todo{Specify which models used, 
bias correction method.}

% ── 3. Methods ───────────────────────────────────────────────────────────────
\section{Methods}
\label{sec:methods}

\subsection{Degree Day Calculation}
\label{sec:methods_dd}

Heating Degree Days (HDD) and Cooling Degree Days (CDD) were calculated 
using the Hungarian standard base temperatures of \SI{15.5}{\celsius} 
and \SI{18.0}{\celsius} respectively:

\begin{align}
  \HDD_d &= \max(T_\mathrm{base,H} - \bar{T}_d,\ 0) \label{eq:hdd} \\
  \CDD_d &= \max(\bar{T}_d - T_\mathrm{base,C},\ 0) \label{eq:cdd}
\end{align}

where $\bar{T}_d$ is the daily mean surface temperature on day $d$.

\subsection{Demand Model}
\label{sec:methods_model}

\todo{Describe XGBoost setup, features, train/test split, 
cross-validation strategy, hyperparameter tuning.}

\subsection{Transmission Loss Model}
\label{sec:methods_losses}

\todo{Describe how grid losses are estimated from ENTSO-E 
generation vs. load delta, and how ambient temperature 
modulates resistance losses.}

% ── 4. Results ───────────────────────────────────────────────────────────────
\section{Results}
\label{sec:results}

\subsection{Historical Temperature–Demand Correlation}
\label{sec:results_historical}

\todo{Figures and analysis here.}

\subsection{Demand Projections to 2050}
\label{sec:results_projections}

\todo{Scenario comparison figures here.}

\subsection{Infrastructure Loss Amplification}
\label{sec:results_losses}

\todo{Loss analysis figures here.}

% ── 5. Discussion ────────────────────────────────────────────────────────────
\section{Discussion}
\label{sec:discussion}

\todo{Interpret results, compare with regional literature, 
discuss limitations and uncertainties.}

% ── 6. Conclusion ────────────────────────────────────────────────────────────
\section{Conclusion}
\label{sec:conclusion}

\todo{Summarise findings and policy implications.}

% ── Acknowledgements ─────────────────────────────────────────────────────────
\section*{Acknowledgements}

ERA5 data were obtained from the Copernicus Climate Change Service. 
ENTSO-E Transparency Platform data are used under their open data policy. 
\todo{Add any other acknowledgements.}

% ── Data Availability ────────────────────────────────────────────────────────
\section*{Data and Code Availability}

All code, processed datasets, and an interactive visualisation are 
publicly available at:\\
\url{https://github.com/Zoomb-o/kisalfold-climate-energy}

% ── References ───────────────────────────────────────────────────────────────
\bibliographystyle{apalike}
\bibliography{references}

% ── Appendix ─────────────────────────────────────────────────────────────────
\appendix

\section{Sounding Station Metadata}
\label{app:soundings}

\todo{Table with sounding station coordinates, elevation, 
collection period, instrument details.}

\section{Model Hyperparameters}
\label{app:hyperparams}

\todo{Full XGBoost hyperparameter table.}

\end{document}
